% =============================================================================
% Philemon-TSH: Tiered-Residency Temporal Subgraph Retrieval on Heterogeneous
%               GPU Memory
%
% Written by filling the DES-LOC paper skeleton
%   (dylanyunlon/Desynced-Low-Communication) with the Philemon-TSH system
%   (dylanyunlon/walpurgis-WTFGG).  Five main body sections, no appendix.
%
% Claude #1 (this file): self-contained preamble, title, abstract, Section 1
%   (Introduction + contributions), and structured stubs for Sections 2-5.
% Claude #2: Section 2  (The Philemon-TSH System)
% Claude #3: Section 3  (Correctness & Complexity Guarantees)
% Claude #4: Section 4  (Experimental Design)
% Claude #5: Section 5  (Evaluation) + short Related Work & Conclusion
% =============================================================================

\documentclass[10pt]{article}

% ---- self-contained bibliography (no external .bib needed) -------------------
\begin{filecontents*}{references.bib}
@inproceedings{paszke2019pytorch,
  title={PyTorch: An Imperative Style, High-Performance Deep Learning Library},
  author={Paszke, Adam and Gross, Sam and Massa, Francisco and others},
  booktitle={NeurIPS}, year={2019}}
@article{shoeybi2019megatron,
  title={Megatron-LM: Training Multi-Billion Parameter Language Models Using Model Parallelism},
  author={Shoeybi, Mohammad and Patwary, Mostofa and Puri, Raul and others},
  journal={arXiv:1909.08053}, year={2019}}
@inproceedings{rasley2020deepspeed,
  title={DeepSpeed: System Optimizations Enable Training Deep Learning Models with Over 100 Billion Parameters},
  author={Rasley, Jeff and Rajbhandari, Samyam and Ruwase, Olatunji and He, Yuxiong},
  booktitle={KDD}, year={2020}}
@inproceedings{abadi2016tensorflow,
  title={TensorFlow: A System for Large-Scale Machine Learning},
  author={Abadi, Mart\'in and Barham, Paul and Chen, Jianmin and others},
  booktitle={OSDI}, year={2016}}
@inproceedings{hamilton2017graphsage,
  title={Inductive Representation Learning on Large Graphs},
  author={Hamilton, William L. and Ying, Rex and Leskovec, Jure},
  booktitle={NeurIPS}, year={2017}}
@article{rossi2020tgn,
  title={Temporal Graph Networks for Deep Learning on Dynamic Graphs},
  author={Rossi, Emanuele and Chamberlain, Ben and Frasca, Fabrizio and others},
  journal={arXiv:2006.10637}, year={2020}}
@inproceedings{erling2015ldbc,
  title={The LDBC Social Network Benchmark: Interactive Workload},
  author={Erling, Orri and Averbuch, Alex and Larriba-Pey, Josep and others},
  booktitle={SIGMOD}, year={2015}}
@article{pugh1990skiplists,
  title={Skip Lists: A Probabilistic Alternative to Balanced Trees},
  author={Pugh, William}, journal={Communications of the ACM},
  volume={33}, number={6}, pages={668--676}, year={1990}}
@book{clrs2009,
  title={Introduction to Algorithms}, edition={3rd},
  author={Cormen, Thomas H. and Leiserson, Charles E. and Rivest, Ronald L. and Stein, Clifford},
  publisher={MIT Press}, year={2009}}
@article{oneil1996lsm,
  title={The Log-Structured Merge-Tree (LSM-Tree)},
  author={O'Neil, Patrick and Cheng, Edward and Gawlick, Dieter and O'Neil, Elizabeth},
  journal={Acta Informatica}, volume={33}, number={4}, pages={351--385}, year={1996}}
@misc{nccl,
  title={{NVIDIA} Collective Communications Library ({NCCL})},
  author={{NVIDIA}}, howpublished={\url{https://github.com/NVIDIA/nccl}}, note={Software}}
@misc{morphgl,
  title={{MorphGL}: Collective Mini-batch Scheduling for {GNN} Training},
  author={{initzhang}}, howpublished={\url{https://github.com/initzhang/MorphGL}}, note={Software}}
@misc{d2stgnn,
  title={{D2STGNN}: Decoupled Dynamic Spatial-Temporal Graph Neural Network},
  author={{GestaltCogTeam}}, howpublished={\url{https://github.com/GestaltCogTeam/D2STGNN}}, note={Software}}
@misc{rapidstore,
  title={{RapidStore}: Concurrent Dynamic Graph Storage with Snapshot Isolation},
  howpublished={Upstream component (see project repository)}, note={Software}}
@misc{temgraph,
  title={{TEM-Graph}: Temporal Interval Index for Dynamic Subgraph Queries},
  howpublished={Upstream component (see project repository)}, note={Software}}
@misc{leveldb,
  title={{LevelDB}}, author={{Google}},
  howpublished={\url{https://github.com/google/leveldb}}, note={Software}}
@misc{cccl,
  title={{CCCL}: {CUDA} Core Compute Libraries ({Thrust}, {CUB}, libcu++)},
  author={{NVIDIA}}, howpublished={\url{https://github.com/NVIDIA/cccl}}, note={Software}}
\end{filecontents*}

\usepackage[margin=1in]{geometry}
\usepackage{amsmath,amssymb,amsfonts,amsthm}
\usepackage{algorithm}
\usepackage{algorithmic}
\usepackage{graphicx}
\usepackage{booktabs}
\usepackage{multirow}
\usepackage{xcolor}
\usepackage[numbers,sort&compress]{natbib}
\usepackage[colorlinks=true,allcolors=blue]{hyperref}

% ---- theorem environments (used by Section 3) --------------------------------
\newtheorem{theorem}{Theorem}
\newtheorem{lemma}{Lemma}
\newtheorem{assumption}{Assumption}
\theoremstyle{definition}
\newtheorem{definition}{Definition}
\newcommand{\todoclaude}[1]{\medskip\noindent\textcolor{red}{\textbf{[TODO #1]}}\medskip}

\title{\textbf{Philemon-TSH}: Tiered-Residency Temporal Subgraph Retrieval\\
       on Heterogeneous GPU Memory}

\author{
  Anonymous Authors\\
  \\
  Walpurgis Project --- Workload-Aware GNN Training on Mixed-Generation GPUs
}
\date{}

\begin{document}
\maketitle

% =============================================================================
\begin{abstract}
Training and serving temporal graph neural networks requires repeatedly
extracting time-windowed subgraphs from graphs whose temporal-edge tables far
exceed the high-bandwidth memory (HBM) of a single accelerator. Existing
temporal stores assume a flat, homogeneous address space: interval indexes such
as TEM-Graph accelerate within-window lookups, and concurrent stores such as
RapidStore provide snapshot isolation, but neither coordinates temporal locality
with the deep, asymmetric memory hierarchy of a mixed-generation cluster
(HBM on an H100, GDDR on an A6000, and host DRAM). Pinning the whole graph in HBM
does not fit; demoting cold data to DRAM without temporal awareness destroys
query latency; and a per-query scan over every temporal partition costs $O(P)$,
which grows without bound under streaming ingestion. We propose \textbf{Philemon-TSH},
a temporal-subgraph engine that assigns \emph{independent residency tiers} to
temporal partitions by access frequency and temporal density---hot, dense
intervals waterfall into HBM, cold, sparse intervals settle in DRAM---and that
prunes partition \emph{selection} to $O(\log P + k)$ with an interval skip list
augmented by a span-maximum of interval endpoints. Our analysis shows that the
inter-partition selection step, not the intra-partition scan, dominates query
cost, and that streaming ingestion requires at least an \emph{incremental}
(segmented, log-structured-merge) index to avoid an $O(N^2\log N)$ rebuild cliff;
we further prove that a sequence lock makes hot-path reads wait-free without
blocking background migration. Across synthetic and LDBC temporal workloads,
Philemon-TSH achieves a $\mathbf{43.9\times}$ faster intra-partition scan than a
linear scan and up to $\mathbf{4.1\times}$ faster partition selection than a full
sweep, at near-zero memory fragmentation under repeated cross-tier migration.
Correctness is validated by over ten million exhaustive cross-checks against
brute-force enumeration and by a $2.1$M-query concurrent run that is race- and
torn-read-free under a thread sanitizer. Unlike single-tier stores, Philemon-TSH
degrades gracefully as graphs outgrow HBM, providing a scalable, race-free,
bandwidth-aware substrate for temporal GNN training on heterogeneous hardware.
\end{abstract}

% =============================================================================
\section{Introduction}
\label{sec:introduction}
% =============================================================================

Temporal graph neural networks operate not on a static adjacency structure but on
a stream of timestamped edges, and almost every training or inference step begins
by materializing a \emph{temporal subgraph}: all edges whose timestamps fall in a
query window $[\ell, h]$, gathered around a batch of seed
vertices~\citep{rossi2020tgn,hamilton2017graphsage}. On modern accelerators this
gather is memory-bound rather than compute-bound, and it is bound by a hierarchy
whose levels differ by more than an order of magnitude in latency: on-package
high-bandwidth memory (HBM) services a temporal-edge lookup in ${\sim}1$\,ns,
the GDDR of a previous-generation board in ${\sim}5$\,ns, and host DRAM behind
PCIe in ${\sim}50$\,ns. Production clusters are rarely homogeneous---a node may
pair a current-generation H100 (large, fast HBM) with one or more A6000-class
boards (ample but slower GDDR), all backed by abundant DRAM---so a temporal store
that targets such hardware must decide \emph{which} time intervals live in
\emph{which} memory, and must revisit that decision as the workload's hot window
slides forward in time.

Prior temporal graph systems do not make this decision. Interval-indexing engines
such as TEM-Graph~\citep{temgraph} build a doubly-linked interval index that
answers contains/contained queries efficiently, but within a single, flat memory;
concurrent dynamic-graph stores such as RapidStore~\citep{rapidstore} provide
snapshot isolation over one address space; and collective mini-batch schedulers
for GNN training such as MorphGL~\citep{morphgl} and decoupled spatial--temporal
models such as D2STGNN~\citep{d2stgnn} assume the working set already fits in
device memory. Meanwhile, the general-purpose caching allocators of deep-learning
frameworks~\citep{paszke2019pytorch,nccl} pool memory \emph{within} a device but
are oblivious to temporal access patterns. Reusing these components on a deep
memory hierarchy poses three problems. \emph{First}, capacity: a billion-edge
temporal-edge table does not fit in HBM, so something must reside in slower tiers.
\emph{Second}, placement: spilling cold data to DRAM with no notion of temporal
locality demotes exactly the intervals a sliding-window workload is about to read,
inflating tail latency. \emph{Third}, and most subtly, \emph{selection}: even with
a perfect intra-partition index, answering a query still requires identifying
which temporal partitions overlap the window, and a linear sweep over all $P$
partitions costs $O(P)$ per query---fixed-cost when $P$ is small, but unbounded
under streaming ingestion, where $P$ grows monotonically as new edges arrive.
Hence we ask:

\begin{center}
\emph{Can temporal partitions be placed across an asymmetric HBM/GDDR/DRAM
hierarchy by their access locality, and selected by an incremental, augmented
index, so that temporal-subgraph retrieval stays sub-millisecond and correct as
the graph outgrows fast memory and edges keep arriving?}
\end{center}

As a result of this inquiry, we propose \textbf{Philemon-TSH}, a temporal-subgraph
engine that composes two orthogonal mechanisms over a common partition abstraction.
The first is a \emph{tiered residency schedule}: a \texttt{TieredAllocator} that
waterfalls each partition across HBM\,$\rightarrow$\,GDDR\,$\rightarrow$\,DRAM, a
\texttt{TemporalBridge} that maps each time interval to a subgraph partition and a
home tier, and a background \texttt{MigrationScheduler} that promotes hot
partitions toward HBM and demotes cold ones toward DRAM. Placement is driven by a
separation of \emph{access timescales} analogous to the separation of optimizer-state
timescales in desynchronized training: dense, recently touched intervals change
tier-membership quickly and belong in fast memory, whereas sparse, stale intervals
are stable and can settle in DRAM. The second mechanism is an \emph{augmented
selection index}: a Pugh skip list~\citep{pugh1990skiplists} ordered by interval
start and augmented, on every forward pointer, with the maximum interval endpoint
over the run that pointer spans---the skip-list analogue of an interval tree's
subtree-maximum augmentation~\citep[\S14.3]{clrs2009}---so that
\texttt{overlaps}$(\ell,h)$ prunes entire runs whose every interval ends before
$\ell$ and returns the $k$ matching partitions in $O(\log P + k)$ rather than
$O(P)$. To keep this index cheap under streaming, the bridge holds it as a
\emph{segmented} structure in the log-structured-merge style~\citep{oneil1996lsm}:
each ingestion flush appends one immutable segment over only the new partitions,
and a threshold-triggered compaction amortizes the merge, sidestepping the
$O(N^2\log N)$ cost of rebuilding a monolithic index on every flush. The residency
schedule and the selection index operate on orthogonal axes---one decides
\emph{where} a partition lives, the other decides \emph{whether} it is touched---so
they compose without interference, and a sequence
lock makes the read path wait-free: queries optimistically read partition metadata
and retry only when a concurrent migration is detected, never blocking and never
forcing the migrating writer to wait on readers.

Following the engineering discipline of the upstream project, each component is
built ``from C, the good example''---a named reference pattern from a
production system is identified and adapted, rather than designed from
scratch. The tiered allocator follows NCCL's allocation dispatch~\citep{nccl};
slab management follows the page-and-bitmask free lists of NCCL and the block
coalescing of the PyTorch caching allocator~\citep{paszke2019pytorch}; the
optimizer-state-swap idiom of DeepSpeed~\citep{rasley2020deepspeed} informs
double-buffered asynchronous migration; and the partition-state grouping echoes
Megatron-LM's grouped allocators~\citep{shoeybi2019megatron}. The upstream
interval semantics come from TEM-Graph~\citep{temgraph} and the concurrent
snapshot protocol from RapidStore~\citep{rapidstore}.

\paragraph{Contributions.}
\begin{enumerate}
  \item \textbf{Access-locality tiered residency for temporal partitions.}
  We introduce a \texttt{TieredAllocator} that waterfalls partitions across an
  asymmetric HBM/GDDR/DRAM hierarchy and a \texttt{TemporalBridge} that maps
  intervals to partitions to home tiers, with a background
  \texttt{MigrationScheduler} that re-tiers partitions by access frequency and
  temporal density. Placement is governed by a separation of access timescales:
  dense, hot intervals are promoted toward HBM while sparse, cold intervals
  settle in DRAM, so the working set of a sliding-window workload stays in fast
  memory without manual pinning.

  \item \textbf{Partition selection in $O(\log P + k)$ via an augmented,
  segmented interval index.} We replace the $O(P)$ per-query partition sweep with
  a skip list ordered by interval start and augmented with a per-pointer
  span-maximum of interval endpoints; \texttt{overlaps}$(\ell,h)$ prunes runs that
  end before $\ell$ and returns the $k$ overlapping partitions in
  $O(\log P + k)$. To stay cheap under streaming, the index is held as a
  log-structured-merge sequence of immutable segments---one segment per flush at
  $O(M\log M)$ for $M$ new partitions, with threshold-triggered compaction---which
  avoids the $O(N^2\log N)$ rebuild cost of a monolithic index.

  \item \textbf{Wait-free reads and fragmentation-free migration under concurrency.}
  We make the hot read path wait-free with a sequence lock: queries read partition
  metadata optimistically and retry only on a detected migration, eliminating the
  writer starvation of a shared-mutex read path. Cross-tier moves route through a
  per-tier slab allocator with $O(1)$ bitmask slot allocation and threshold-driven
  compaction, keeping memory fragmentation near zero under repeated
  promotion/demotion cycles, and through double-buffered asynchronous copies that
  overlap migration with queries.

  \item \textbf{Validated correctness and measured speedups on heterogeneous hardware.}
  We achieve a $43.9\times$ faster intra-partition scan than a linear scan
  (a narrow window drops from $232.6\,\mu$s to $5.3\,\mu$s) and up to $4.1\times$
  faster partition selection than a full sweep, and we validate correctness by
  over ten million exhaustive cross-checks against brute-force enumeration,
  a twenty-thousand-query three-way check of the segmented index against a
  whole-list index and brute force, and a $2.1$M-query concurrent run that a
  thread sanitizer confirms is free of data races and torn reads.
\end{enumerate}

% =============================================================================
\section{The Philemon-TSH System}
\label{sec:system}
% =============================================================================
% --- Claude #2 ---  (template analogue: DES-LOC Section 2, "DES-LOC Algorithm")

Philemon-TSH organizes a temporal graph as a set of \emph{temporal partitions}.
Each partition holds a contiguous, sorted run of timestamped edges together with
a per-partition interval index and a residency tier, and carries an interval
$[\mathrm{ts\_lo},\mathrm{ts\_hi}]$ spanning its earliest start and latest end.
Over this abstraction the system composes two mechanisms on orthogonal axes:
a \emph{residency schedule} that decides \emph{where} a partition lives in the
HBM/GDDR/DRAM hierarchy (\S\ref{sec:timescales}), and a \emph{selection index}
that decides \emph{whether} a partition is touched by a query
(\S\ref{sec:architecture}). Following the project's ``from \texttt{C}, the good
example'' discipline, each component adapts a named pattern from a production
system---RapidStore's \texttt{snapshot\_edges} traversal abstraction is the
seed~\citep{rapidstore}---rather than being designed from scratch.

\subsection{Separation of Access Timescales}
\label{sec:timescales}

A temporal partition $P$ is touched by a query only when its interval meets the
query window, so the rate at which $P$ enters and leaves the active working
set---its \emph{access timescale}---determines how valuable a fast residency tier
is for it. Two signals govern that timescale. The first is \emph{temporal
density}
\[
  \rho(P) \;=\; \frac{\text{edge\_count}(P)}{\mathrm{ts\_hi}(P) - \mathrm{ts\_lo}(P) + 1}.
\]
A dense partition packs many edges into a narrow time range, so a sliding window
dwells on it briefly but intensely; a sparse partition spreads few edges over a
wide range and is read faintly but for a long stretch of the timeline. The second
is \emph{access frequency}: every read increments a partition's lock-free
\texttt{access\_count} and stamps \texttt{last\_access\_ns}, recording how often
and how recently it was used. Dense, recently touched partitions form the
fast-churning hot set and must be cheap to reach (HBM); sparse, stale partitions
are stable and tolerate the higher latency of DRAM because they are read rarely.
This is the residency analogue of the separation of optimizer-state timescales in
desynchronized training, where fast-moving parameters are averaged often and
slow-moving second moments rarely: here it is the \emph{residency tier}, not the
synchronization period, that differs along the partition's access timescale.

Two places act on this separation. \textbf{Flush-time partitioning} makes density
explicit. When the bridge flushes its ingest buffer it computes the average
density $\bar\rho$ over the batch and walks the start-sorted edges in a
look-ahead window; a region whose local density exceeds $2\bar\rho$ is cut into
\emph{smaller} partitions sized to fit HBM, while a region below $0.5\bar\rho$ is
coalesced into a \emph{larger} partition destined for DRAM. This density-adaptive
granularity follows TEM-Graph's density-aware index construction~\citep{temgraph}.
\textbf{Runtime migration} acts on frequency: a background scheduler re-tiers
partitions by a hotness score $h(P)$ combining \texttt{access\_count} and
recency, promoting hot partitions toward HBM and demoting cold ones toward DRAM,
subject to per-tier capacity budgets (on the target server, $80$\,GB HBM on the
H100, $48$\,GB GDDR on the A6000, and host DRAM behind PCIe). Because the
scheduler runs deferred and batched---the grouped, deferred-hook pattern of
Megatron-LM's data-parallel reducer~\citep{shoeybi2019megatron}---it keeps
migration off the query critical path.

\subsection{The Philemon-TSH Architecture}
\label{sec:architecture}

\paragraph{Tiered allocator.}
\texttt{Allocate}$(n,\text{preferred})$ tries to reserve $n$ bytes in the
preferred tier and, if its budget is exhausted, waterfalls
HBM\,$\rightarrow$\,GDDR\,$\rightarrow$\,DRAM, each reservation a lock-free
compare-and-swap on the tier's used-bytes counter. This dispatch-by-availability
mirrors NCCL's \texttt{ncclMemAlloc}, which falls through
\texttt{cuMemCreate}/\texttt{cuMemMap}/\texttt{malloc} by
capability~\citep{nccl}; in production, HBM and GDDR map to device allocations on
the H100 and A6000 and DRAM to pinned host memory, while the CPU-only development
harness simulates all three from DRAM. Per-allocation metadata---size, home tier,
base pointer, and atomic \texttt{access\_count} and \texttt{last\_access\_ns}---lives
in a registry guarded by a shared mutex: structural mutations
(\texttt{allocate}/\texttt{deallocate}/\texttt{migrate}) take it exclusively,
ordinary reads share it, and \texttt{touch} takes no lock at all, being a single
relaxed atomic fetch-add following NCCL's atomic refcount and CCCL's lock-free
shared block pointer~\citep{cccl}. Each tier routes small allocations
($\le 512$\,KB) through a per-tier slab allocator whose pages carry a $64$-bit
free mask; a slot is claimed in $O(1)$ via \texttt{\_\_builtin\_ctzll} on the mask
(NCCL's page-and-bitmask pooling), freed slots coalesce, and \texttt{compact}
returns wholly empty pages to the OS---the block-coalescing idiom of the PyTorch
caching allocator~\citep{paszke2019pytorch} over a TensorFlow-Arena
bump-pointer fast path~\citep{abadi2016tensorflow}. Fragmentation thus stays near
zero across repeated \texttt{migrate} cycles.

\paragraph{Temporal bridge.}
The bridge realizes interval\,$\rightarrow$\,partition\,$\rightarrow$\,tier.
\texttt{Ingest} appends edges to a buffer; \texttt{Flush} sorts the buffer by
$(\mathrm{ts\_start},\mathrm{ts\_end})$, cuts it into the density-adaptive
partitions of \S\ref{sec:timescales}, allocates each through the tiered
allocator, and builds a per-partition \emph{interval index}---two sorted views of
the partition's edges, one by start and one by end, the dual-ordered-set pattern
of the PyTorch \texttt{BlockPool}~\citep{paszke2019pytorch}---so an in-partition
contains/contained/overlaps query is answered in $O(\log N + k)$ by binary
search rather than an $O(N)$ linear scan. The in-partition scan itself is
LevelDB's two-level \texttt{Seek}~\citep{leveldb} over Thrust's
\texttt{lower\_bound}~\citep{cccl}: a binary search to the first edge with
$\mathrm{ts\_start}\ge\ell$, then a forward walk with early termination once
$\mathrm{ts\_start}>h$.

\paragraph{Selection index.}
Answering a query first requires finding \emph{which} partitions meet the window.
A linear sweep is $O(P)$ and grows without bound under streaming ingestion, so the
bridge holds an augmented interval skip list~\citep{pugh1990skiplists}: nodes
ordered by $\mathrm{ts\_lo}$, with every forward pointer carrying the maximum
$\mathrm{ts\_hi}$ over the run of nodes it spans---the skip-list counterpart of an
interval tree's subtree-maximum augmentation~\citep[\S14.3]{clrs2009}.
\texttt{Overlaps}$(\ell,h)$ walks the bottom level but, at higher levels, skips
any run whose span-maximum is below $\ell$ (no member can reach the window) and
stops once a node's $\mathrm{ts\_lo}$ exceeds $h$, returning the $k$ overlapping
partitions in $O(\log P + k)$. To keep this cheap under streaming, the index is
\emph{segmented} in the log-structured-merge style~\citep{oneil1996lsm}: each
flush builds one immutable segment over only its $M$ new partitions in
$O(M\log M)$, queries fan out across segments, and a compaction merges them into a
single segment once their count crosses a threshold of eight. The index is
therefore never rebuilt wholesale, sidestepping the $O(N^2\log N)$ cost of
re-augmenting a monolithic structure on every flush.

\paragraph{Read path and migration.}
The query read path is wait-free. Partition metadata is protected by a sequence
lock: a query samples the sequence number (even $\Rightarrow$ no writer active),
selects and scans the matching partitions, and retries only if a concurrent
migration left the sequence odd or changed; it never blocks, and the migrating
writer never waits on readers, eliminating the writer starvation of the
shared-mutex read path it replaces (the Linux-kernel seqlock and NCCL sequence
numbers). Cross-tier moves are issued by the background scheduler and executed
double-buffered and asynchronously through a scope-safe handle---an RAII
\texttt{TierPtr} in the manner of PyTorch's \texttt{CUDAStreamGuard}~\citep{paszke2019pytorch}---so
a migrating copy overlaps in-flight queries, with an event fence enforcing
read-after-write ordering. Algorithm~\ref{alg:philemon} summarizes the three
paths.

\begin{algorithm}[t]
\caption{Philemon-TSH ingest, query, and background migration}
\label{alg:philemon}
\begin{algorithmic}[1]
\STATE \textbf{Shared state:} partitions $\mathcal{P}$; segmented selection index $\mathcal{I}$; tiered allocator $\mathcal{A}$; sequence lock $\sigma$
\STATE
\STATE \textbf{procedure} \textsc{Flush}( ) \COMMENT{fires when the ingest buffer $B$ overflows}
\STATE \quad sort $B$ by $(\mathrm{ts\_start},\mathrm{ts\_end})$; \quad $\bar\rho \gets |B| / (\max\,\mathrm{ts\_start} - \min\,\mathrm{ts\_start})$
\FOR{each density-adaptive partition $P=[\mathrm{ts\_lo},\mathrm{ts\_hi}]$ cut from $B$ \quad ($\rho>2\bar\rho$: smaller; $\rho<0.5\bar\rho$: larger)}
  \STATE $P.\mathrm{alloc} \gets \mathcal{A}.\textsc{Allocate}(\mathrm{size}(P),\ \mathrm{HBM})$ \COMMENT{waterfall HBM\,$\to$\,GDDR\,$\to$\,DRAM}
  \STATE build the dual-sorted interval index for $P$ \COMMENT{by start, by end}
\ENDFOR
\STATE append one immutable segment over the new partitions to $\mathcal{I}$ \COMMENT{$O(M\log M)$}
\IF{$\mathcal{I}$ holds $> 8$ segments}
  \STATE compact $\mathcal{I}$ into one segment
\ENDIF
\STATE
\STATE \textbf{procedure} \textsc{Query}($\ell,h$, cb)
\REPEAT
  \STATE $s \gets \sigma.\textsc{ReadBegin}()$ \COMMENT{spin while a writer is active}
  \STATE $\mathrm{slots} \gets \mathcal{I}.\textsc{Overlaps}(\ell,h)$ \COMMENT{span-max pruning, $O(\log P + k)$}
  \FORALL{$\mathrm{slot} \in \mathrm{slots}$}
    \STATE $P \gets \mathcal{P}[\mathrm{slot}]$; \quad $i \gets \textsc{LowerBound}(P.\mathrm{edges},\ \ell)$ \COMMENT{$O(\log N)$}
    \STATE emit $P$'s edges in $[\ell,h]$ from $i$; \quad $\mathcal{A}.\textsc{Touch}(P.\mathrm{alloc})$ \COMMENT{lock-free}
  \ENDFOR
\UNTIL{$\lnot\,\sigma.\textsc{ReadRetry}(s)$} \COMMENT{retry only on a torn read}
\STATE
\STATE \textbf{procedure} \textsc{MigrateSweep}( ) \COMMENT{background thread}
\FORALL{$P \in \mathcal{P}$}
  \STATE $h(P) \gets \mathrm{hotness}(P.\mathrm{access\_count},\ P.\mathrm{last\_access})$
  \IF{$h(P)$ high $\,\land\, P \notin \mathrm{HBM} \,\land\,$ HBM budget available}
    \STATE $\sigma.\textsc{WriteLock}()$;\ \ $\mathcal{A}.\textsc{Migrate}(P,\ \mathrm{HBM})$;\ \ $\sigma.\textsc{WriteUnlock}()$
  \ELSIF{$h(P)$ low}
    \STATE $\sigma.\textsc{WriteLock}()$;\ \ $\mathcal{A}.\textsc{Migrate}(P,\ \text{next colder tier})$;\ \ $\sigma.\textsc{WriteUnlock}()$
  \ENDIF
\ENDFOR
\end{algorithmic}
\end{algorithm}

% =============================================================================
\section{Correctness and Complexity Guarantees}
\label{sec:guarantees}
% =============================================================================
% --- Claude #3 fills this section ---
% Template analogue: DES-LOC Section 3 (Convergence Guarantees).
% Required structure:
%   - Assumptions/Invariants (analogue of Assumptions 1-3): partitions sorted by
%     ts_lo on flush; closed intervals; span-max >= every member ts_hi.
%   - Theorem (selection cost & soundness): augmented skip-list overlaps(lo,hi)
%     is O(log P + k); the span_max < lo prune is sound (every interval in a
%     skipped run ends before lo). (analogue of Theorem 1.)
%   - Amortization analysis (analogue of the psi-factor boundary analysis):
%     per-flush O(M log M) segment build; compaction at kCompactThreshold = 8
%     gives amortized build, no O(N^2 log N) cliff.
%   - SeqLock read-consistency statement (torn-read detection).
%   - Validation harness (analogue of the Fig. 1 toy problem): 10M exhaustive +
%     32k random/adversarial + 20k three-way cross-checks; 900/900 runtime
%     samples by alloc_id; 2.1M-query TSan run, no race.
% Grounding: src/core/partition_skiplist.hpp, src/bench/skiplist_selftest.cpp,
%   src/core/seqlock.hpp; REVIEW_M013_M014.md.
\todoclaude{Claude \#3: Correctness and Complexity Guarantees (\S\ref{sec:guarantees}).}

% =============================================================================
\section{Experimental Design}
\label{sec:expdesign}
% =============================================================================
% --- Claude #4 fills this section ---
% Template analogue: DES-LOC Section 4 (Experimental Design) + 4.1 Setup.
% Required structure:
%   - Research questions RQ1-RQ6 (e.g. RQ1: does the augmented index predict the
%     measured selection cost? RQ2: how does tier placement affect latency?
%     RQ3: scan/selection speedup vs linear? RQ4: scaling to 100M edges?
%     RQ5: behaviour under streaming ingestion + compaction? RQ6: concurrency /
%     migration overlap?).
%   - \subsection{Experimental Setup}: workloads (synthetic 1M-edge temporal
%     graph + LDBC SNB SF-1/SF-10/SF-100), tier configuration (HBM 512 MB /
%     GDDR 1024 MB / DRAM 2048 MB in the dev harness), baselines (single-tier
%     all-HBM, all-DRAM, linear-sweep oracle, TEM-Graph-only, RapidStore-only),
%     metrics (query latency by window width, QPS, per-tier memory utilization,
%     migration cost, scan & selection speedup), hardware (H100 + A6000 + DRAM;
%     sm_86 + compute_80 PTX), and the 2000-points x 3-seeds data protocol.
% Grounding: Makefile, build_hetero.sh, run_hetero.sh, DEVELOPMENT_PLAN.md,
%   philemon_*_2000.json headers.
\todoclaude{Claude \#4: Experimental Design (\S\ref{sec:expdesign}).}

% =============================================================================
\section{Evaluation}
\label{sec:evaluation}
% =============================================================================
% --- Claude #5 fills this section (+ short Related Work & Conclusion below) ---
% Template analogue: DES-LOC Section 5 (Evaluation), one subsection per RQ.
% Use the philemon_*_2000.json curves (2000 pts x 3 seeds): query_latency,
% qps, memory_util, migration_cost, interval_index, partition_index,
% async_migration, query_under_mig. Report means +/- std of the final window.
% Key results to land: intra-partition scan 43.9x (narrow 232.6->5.3 us;
% medium 108.5 us; wide 608 us; full 1230 us at ~1.23 ns/edge); partition
% selection at P=8000 narrow 4.1x / medium 1.75x / wide 0.85x, pure selection
% 3.7 us indexed vs 8.0 us linear; near-zero fragmentation; correctness +
% TSan results cross-referenced to Section 3.
% Grounding: REVIEW_M005_M006.md, REVIEW_M013_M014.md, philemon_*_2000.json.
\todoclaude{Claude \#5: Evaluation (\S\ref{sec:evaluation}).}

% =============================================================================
% --- Claude #5 also appends, after Evaluation (short, DES-LOC 6-7 length): ---
%   \section{Related Work}     temporal graph stores (TEM-Graph, RapidStore),
%       tiered/heterogeneous memory & caching allocators (PyTorch, NCCL,
%       DeepSpeed swap), interval/skip-list indexing, LSM indexing, GNN systems
%       (MorphGL, D2STGNN).
%   \section{Conclusion}       restate the tiered-residency + augmented-index
%       contribution and the measured speedups; one paragraph.
% =============================================================================

\bibliographystyle{plainnat}
\bibliography{references}

\end{document}
